\chapter{Дүгнэлт}

Энэхүү дипломын ажлын хүрээнд annotation болон reflection механизмд суурилсан dependency injection framework-ийг давхаргат архитектурын зарчмын дагуу Java SE 17 орчинд хэрэгжүүллээ. Annotation давхарга, scanner давхарга, container давхарга, injector давхарга болон exception давхарга гэсэн үндсэн бүрэлдэхүүн хэсгүүд нь тус бүр өөрийн үүргийг тодорхой хязгаарын дотор гүйцэтгэж, хоорондын хамаарлыг интерфэйс болон container-оор зохицуулсан. Гуравдагч талын нэмэлт сан ашиглахгүйгээр зөвхөн JDK стандарт сангийн Reflection API болон annotation inspection механизмыг ашигласан нь framework-ийн дотоод ажиллагааг ил тод харуулах судалгааны зорилгод нийцэв.

Хэрэгжүүлэлтийн үр дүнд \texttt{@Component}, \texttt{@Autowired}, \texttt{@Qualifier}, \texttt{@Scope} annotation-уудыг дэмжих, singleton болон prototype scope удирдах, bean-ийг нэр болон төрлөөр авах, dependency injection хийх, тойрог хамаарал илрүүлэх, тодорхой exception шидэх зэрэг системийн шаардлагын бүлэгт тодорхойлсон үндсэн боломжуудыг хангаж чадсан. Энэ нь 2-р бүлгийн шаардлага, 3-р бүлгийн зохиомж, 4-р бүлгийн хэрэгжүүлэлтийн хооронд бодит уялдаа тогтсоныг харуулж байна.

Мөн нэмэлт хэрэгжүүлэлтийн A1 ажлын хүрээнд одоо байгаа хэрэгжилтийн чанарын баталгаажуулалтыг өргөжүүлж, fail-fast exception-уудын мэдээллийн чанарыг сайжруулан, \texttt{unit} болон \texttt{integration} түвшний туршилтын санг цэгцэлсэн. Үүний дүнд нийт 10 тест класс, 47 тестийн тохиолдлоор framework-ийн гол болон захын урсгалуудыг баталгаажуулж, \texttt{mvn test} гүйцэтгэлээр амжилттай шалгав. Иймээс энэхүү framework нь зөвхөн ажиллаж буй туршилтын загвар бус, харин шаардлага, зохиомж, туршилтын нотолгоогоор бэхжсэн хэрэгжүүлэлт болсон гэж дүгнэж байна.

Цаашид нэмэлт хэрэгжүүлэлтийн A2, A3, A4 ажлуудын хүрээнд constructor injection, нэмэлт annotation, dependency tree, log бичих систем зэрэг боломжуудыг шат дараатайгаар өргөтгөх боломжтой. Харин эдгээр дараагийн ажлын суурь нь A1-ээр бий болсон тестийн хамгаалалт, exception-ийн ойлгомжтой байдал, давхарга тус бүрийн хариуцлагыг нотолсон чанарын баталгаажуулалт юм.
